% !TeX spellcheck = ru_RU
% !TEX root = vkr.tex

\section{Программная реализация} \label{sec:implementation}

Для проведения экспериментов по исследованию влияния дообучения большой языковой модели на качество классификации инструментов был разработан программный комплекс \texttt{ml\_benchmark}. Система реализована на языке Python~3.11 и обеспечивает полный цикл экспериментальной работы: от генерации синтетических данных до статистического анализа результатов.

\subsection{Архитектура системы}

Программный комплекс построен по принципу модульной архитектуры с чётким разделением ответственности между слоями. Верхний уровень образует интерактивный CLI-интерфейс, реализованный с помощью библиотеки Rich и предоставляющий меню из 11 команд, каждая из которых соответствует отдельному этапу экспериментального процесса. Средний уровень составляет бизнес-логика, сосредоточенная в модуле \texttt{core/}: генерация данных, инференс моделей, вычисление метрик, статистический анализ, визуализация и формирование отчётов. Нижний уровень включает подсистемы хранения (\texttt{storage/}), конфигурации (\texttt{config/}) и LLM-клиентов (\texttt{clients/}).

Взаимодействие между слоями организовано через паттерн внедрения зависимостей (Dependency Injection). Главный цикл приложения, реализованный в классе \texttt{App} (модуль \texttt{cli.py}), при инициализации создаёт четыре ключевых объекта: \texttt{settings} (конфигурация из YAML-файлов и переменных окружения), \texttt{artifacts} (менеджер файловых артефактов), \texttt{registry} (реестр моделей) и \texttt{saves} (менеджер сохранений). Каждый обработчик команды получает эти объекты в качестве параметров, что обеспечивает слабую связанность компонентов и упрощает тестирование.

Для строгой типизации всех структур данных используются Pydantic-модели, определённые в пакете \texttt{models/}. Они описывают схемы датасетов, экспериментов, метрик, отчётов групповой оценки и итоговых отчётов. Валидация данных на уровне моделей позволяет обнаруживать ошибки на ранних стадиях и гарантирует целостность данных при их передаче между компонентами.

Подсистема LLM-клиентов реализована по принципу plugin-архитектуры. Абстрактный базовый класс \texttt{BaseClient} определяет интерфейс взаимодействия с языковыми моделями, включая типы \texttt{Message} и \texttt{ChatResponse}. Конкретные реализации, такие как \texttt{YandexCloudClient}, наследуют этот интерфейс и инкапсулируют специфику работы с API конкретных провайдеров. Такой подход позволяет подключать новых провайдеров без модификации существующего кода.

\subsection{Подсистема статистического анализа}

Центральным компонентом подсистемы статистического анализа является класс \texttt{GroupEvaluator}, реализованный в модуле \texttt{core/group\_evaluator.py}. Данный класс обеспечивает статистически корректное сравнение нескольких моделей на одном и том же наборе тестовых примеров. Алгоритм его работы включает пять основных этапов.

На первом этапе выполняется построение бинарной матрицы результатов размерности $n_{\text{samples}} \times n_{\text{models}}$, где каждый элемент $b_{ij} \in \{0, 1\}$ отражает корректность предсказания $j$-й модели на $i$-м примере. Классификация предсказания как верного (TP, TN) или ошибочного (FP, FN) осуществляется модулем \texttt{Evaluator}, который выполняет парсинг вызовов инструментов из XML-ответов модели и сопоставляет их с эталонными значениями.

На втором этапе для каждой модели вычисляются доверительные интервалы метрик с помощью бутстрапа. Используется $B = 1000$ итераций ресемплирования с возвращением, что обеспечивает непараметрическую оценку распределения метрик без предположений о их аналитическом виде.

На третьем этапе выполняется проверка глобальной гипотезы $H_0$ о равенстве долей правильных ответов всех моделей с помощью критерия Кохрена~$Q$ (реализация из библиотеки \texttt{statsmodels}). Данный критерий является обобщением критерия Макнемара на случай $k > 2$ связанных выборок и применим к бинарным данным при условии, что все модели оцениваются на одних и тех же примерах.

Четвёртый этап выполняется только при отклонении глобальной гипотезы и включает попарные сравнения моделей с помощью критерия Макнемара. Для контроля ошибки первого рода при множественных сравнениях применяется поправка Бонферрони, корректирующая уровень значимости пропорционально числу пар.

На пятом этапе проводится анализ ошибок: для каждой модели формируется перечень примеров с ошибочной классификацией, что позволяет выявить систематические паттерны ошибок и понять, на каких типах запросов модель испытывает затруднения.

\subsection{Подсистема инференса}

Подсистема инференса реализована в классе \texttt{InferenceRunner} (модуль \texttt{core/runner.py}) и обеспечивает запуск языковых моделей на наборе тестовых примеров с учётом ограничений внешних API.

Для обеспечения устойчивости к сбоям реализован механизм чекпоинтов: после обработки каждого примера промежуточные результаты сохраняются на диск. При повторном запуске система автоматически обнаруживает последний чекпоинт и продолжает выполнение с того примера, на котором произошло прерывание. Это критически важно при работе с большими датасетами, поскольку один полный прогон может занимать значительное время.

Для соблюдения ограничений частоты запросов (rate limits) внешних API между последовательными вызовами модели выдерживается пауза в 0.5~секунды. Кроме того, реализована стратегия повторных попыток (retry) при получении временных ошибок от API, что повышает надёжность сбора данных.

Парсинг ответов модели выполняется классом \texttt{ToolCallParser} (модуль \texttt{core/evaluator.py}), который извлекает вызовы инструментов из XML-разметки ответа. Извлечённые вызовы сопоставляются с эталонными, после чего класс \texttt{MetricsCollector} подсчитывает количество истинно положительных (TP), ложноположительных (FP), истинно отрицательных (TN) и ложноотрицательных (FN) результатов и на их основе вычисляет итоговые метрики классификации.

\subsection{Подсистема генерации данных}

Генерация синтетических тестовых данных выполняется двумя компонентами. Класс \texttt{Generator} (модуль \texttt{core/generator.py}) формирует синтетические пользовательские запросы с помощью вызовов LLM, обеспечивая контроль баланса классов в итоговом датасете. Класс \texttt{Humanizer} (модуль \texttt{core/humanizer.py}) выполняет <<очеловечивание>> сгенерированных запросов~--- повторно обрабатывает синтетические примеры через LLM, чтобы придать им стилистическое разнообразие, характерное для реальных пользовательских формулировок. Полученные данные конвертируются в формат JSONL классом \texttt{DatasetConverter} для последующего использования при дообучении моделей.

\subsection{Обеспечение качества кода}

Качество программного комплекса поддерживается комплексом автоматизированных инструментов. Юнит-тестирование осуществляется с помощью фреймворка \texttt{pytest}; покрытие кода тестами составляет порядка 85\%. Статический анализ кода выполняется линтером \texttt{ruff}, а проверка типов~--- инструментом \texttt{mypy}. Совместное использование статической типизации \texttt{mypy} и валидации данных через Pydantic-модели во время исполнения обеспечивает двухуровневый контроль корректности: ошибки несовместимости типов обнаруживаются как на этапе разработки, так и при передаче данных между компонентами в процессе выполнения.

\subsection{Технологический стек}

Перечень основных библиотек и инструментов, используемых в проекте, приведён в таблице~\ref{tab:tech_stack}.

\begin{table}[H]
\centering
\caption{Технологический стек проекта}
\label{tab:tech_stack}
\begin{tabular}{L{0.28\textwidth} L{0.65\textwidth}}
\toprule
\textbf{Библиотека} & \textbf{Назначение} \\
\midrule
Python 3.11       & Язык реализации \\
Pydantic          & Описание и валидация схем данных \\
NumPy             & Числовые операции, работа с матрицами результатов \\
statsmodels       & Критерий Кохрена $Q$, критерий Макнемара \\
matplotlib        & Построение графиков и визуализаций \\
Rich              & Интерактивный CLI-интерфейс \\
PyYAML            & Чтение конфигурационных файлов \\
pytest            & Юнит-тестирование \\
ruff              & Линтинг и форматирование кода \\
mypy              & Статическая проверка типов \\
\bottomrule
\end{tabular}
\end{table}
